% =====================================================================
%  MMPP_- リポジトリ 実装仕様書
%  コンパイル: lualatex docs/model_and_methods.tex  (luatexja を使用)
%  upLaTeX を使う場合は documentclass を
%    \documentclass[a4paper,11pt,uplatex,dvipdfmx]{jsarticle}
%  に変更し, \usepackage{luatexja} 行を削除してください.
% =====================================================================
\documentclass[a4paper,11pt]{article}
\usepackage{luatexja}
\usepackage{amsmath,amssymb}
\usepackage{booktabs}
\usepackage{longtable}
\usepackage[margin=25mm]{geometry}
\usepackage{hyperref}
\hypersetup{colorlinks=true,linkcolor=black,citecolor=black,urlcolor=blue}

\newcommand{\code}[1]{\texttt{#1}}
\newcommand{\off}{\mathrm{off}}
\newcommand{\eff}{\mathrm{eff}}
\newcommand{\blk}{\mathrm{block}}

\title{MMPP/M/$c$/SET-BATCH/Delayoff モデルおよび\\
その Predictive 拡張の実装仕様\\
\large --- モデル定義・数値解析手法・数値実験手順の記述 ---}
\author{本リポジトリ (\code{MMPP\_-}) の実装に基づく記述}
\date{\today}

\begin{document}
\maketitle

\begin{abstract}
本書は本リポジトリに実装されている 2 つの待ち行列モデル,
すなわち (i) ベースモデル \textbf{Delayoff モデル}
(\code{mmpp/}, \code{mmpp\_sim/}) と
(ii) その拡張である \textbf{Predictive モデル}
(\code{mmpp\_predictive/}, \code{mmpp\_predictive\_sim/}) について,
実装されているとおりのモデル定義・性能指標・定常分布の数値解法・
離散事象シミュレーション・数値実験の設計を記述したものである。
数式はすべて実際のソースコードから起こしており, 実装と異なる理想化や
一般化は意図的に行っていない。実装上の非自明な選択・近似・仕様の細部は
本文中および第 \ref{sec:checklist} 節のチェックリストに明示した。
\end{abstract}

\tableofcontents

% =====================================================================
\section{記号と対象コード}
% =====================================================================

\subsection{対象ディレクトリ}

\begin{center}
\begin{tabular}{ll}
\toprule
ディレクトリ & 内容 \\
\midrule
\code{mmpp/}                 & ベースモデルの理論解析 (生成行列・定常分布・性能指標) \\
\code{mmpp\_sim/}            & ベースモデルの離散事象シミュレーション (DES) \\
\code{mmpp\_predictive/}     & Predictive モデルの理論解析 \\
\code{mmpp\_predictive\_sim/}& Predictive モデルの DES \\
\code{scripts/}              & 数値実験スクリプト (実験 0--6, 比較・検証) \\
\code{tests/}                & pytest による単体・結合テスト \\
\code{results/}              & 実験結果 CSV (GTH, 正本) \\
\code{results\_gth/}          & \code{results/} と同一内容の重複 (検証レポートを含む) \\
\code{figures/}              & 出力図 (PNG/PDF, 正本) \\
\code{figures\_gth/}          & 引き算版の式で描いた旧図 (論文には使わない) \\
\bottomrule
\end{tabular}
\end{center}

\code{mmpp\_sim} / \code{mmpp\_predictive\_sim} は理論実装の検証用であり,
補助関数 (\code{busy\_servers} 等) を意図的に独立再実装している
(\code{mmpp.generator} を import しない)。共有するのはパラメータクラスのみである。

\subsection{記号}

\begin{center}
\begin{tabular}{ll}
\toprule
記号 & 意味 \\
\midrule
$c$      & サーバー総数 ($c>0$) \\
$K$      & バッファ容量 (系内ジョブ数の上限, $K \ge b$) \\
$b$      & バッチサイズ ($b>0$, フルバッチ方式) \\
$\mu$    & バッチサービス完了率 ($>0$) \\
$\alpha$ & セットアップ完了率 ($>0$) \\
$\beta$  & Delayoff タイムアウト率 ($>0$) \\
$D_M$    & MMPP の位相数 \\
$C_0,C_1$& MMPP の内部遷移行列 / 到着遷移行列 ($D_M\times D_M$) \\
$i$      & アクティブサーバー数 (Busy または Idle=Delayoff 計時中), $0\le i\le c$ \\
$s$      & セットアップ中サーバー数 (Predictive モデルのみ明示的状態), $0\le s\le c-i$ \\
$j$      & 系内ジョブ数, $0\le j\le K$ \\
$F$      & MMPP 位相, $0\le F<D_M$ \\
$n_{\mathrm{target}}$ & 事前セットアップの目標稼働台数 (Predictive) \\
$\gamma$ & Delayoff 加速係数 (Predictive, $\gamma\ge1$) \\
\bottomrule
\end{tabular}
\end{center}

すべての事象時間は指数分布に従い, 系は連続時間マルコフ連鎖 (CTMC) である。

% =====================================================================
\section{共通の到着過程: MMPP}
% =====================================================================

到着は Markov Modulated Poisson Process (MMPP) である。位相過程は
生成行列 $C_0+C_1$ をもつ $D_M$ 状態 CTMC で, 位相 $F$ にいる間
到着率 $\lambda_F$ のポアソン到着が発生する。実装は $(C_0,C_1)$ の組で与える
(\code{mmpp/model.py})。

\paragraph{整合性条件と検証}
\code{ModelParameters.\_validate} が以下を検査する。
\begin{itemize}
  \item $C_0$ の非対角成分 $\ge 0$, 対角成分 $\le 0$, $C_1 \ge 0$ (許容誤差 $10^{-12}$)。
  \item $(C_0+C_1)\mathbf{e}=\mathbf{0}$ (許容誤差 $10^{-10}$)。
  \item $C_0+C_1$ の非零パターンを有向グラフとみて強連結成分が 1 つであること
        (位相過程の既約性)。これが $Q$ 全体の既約性の十分条件であり,
        定常分布の一意性を保証する。
  \item $\lfloor K/b\rfloor < c$ のとき, $i>\lfloor K/b\rfloor$ の状態が
        セットアップ経路から到達不能になるため \code{UserWarning} を出す
        (数学的には問題ないがパラメータ設定ミスの検出用)。
\end{itemize}

\paragraph{派生量}
\begin{align}
  \lambda_F &= (C_1\mathbf{e})_F, \\
  \boldsymbol{\varpi}(C_0+C_1) &= \mathbf{0},\quad \boldsymbol{\varpi}\mathbf{e}=1
  \qquad (\text{位相定常分布}), \\
  \bar{\lambda} &= \boldsymbol{\varpi}\,\boldsymbol{\lambda}
  \qquad (\text{平均到着率}).
\end{align}
$\boldsymbol{\varpi}$ は $(C_0+C_1)^{\mathsf T}$ の最終行を $\mathbf{1}^{\mathsf T}$ に
置き換えた線形系を \code{numpy.linalg.solve} で解いて得る。

% =====================================================================
\section{ベースモデル (Delayoff モデル)}
\label{sec:base}
% =====================================================================

\subsection{状態空間とインデックス付け}

状態は 3 つ組 $(i,j,F)$ である。
\begin{equation}
  \mathcal{S}=\{(i,j,F): 0\le i\le c,\ 0\le j\le K,\ 0\le F<D_M\}.
\end{equation}
$D_S=c+1$, $D=D_S D_M$, $N=D(K+1)$ とおき, $j$-major の順序
\begin{equation}
  \mathrm{idx}(i,j,F)=jD+iD_M+F
  \label{eq:base-index}
\end{equation}
を用いる (\code{mmpp/state\_space.py})。この順序でレベル $j$ の状態が
連続したブロックになり, $Q$ がブロック帯構造を持つ。

\textbf{重要}: ベースモデルではセットアップ中サーバー数は状態変数ではなく,
$(i,j)$ から決定論的に定まる関数である (次節)。

\subsection{サーバー内訳の補助関数}

\code{mmpp/generator.py} で定義される。
\begin{align}
  B(i,j) &= \min\!\left(i,\ \left\lfloor j/b \right\rfloor\right)
  &&\text{(処理中=Busy)} \label{eq:B}\\
  I(i,j) &= i - B(i,j)
  &&\text{(アイドル; Delayoff 計時中)} \label{eq:I}\\
  S(i,j) &= \max\!\left(0,\ \min\!\left(c,\left\lfloor j/b\right\rfloor\right)-i\right)
  &&\text{(セットアップ中)} \label{eq:S}\\
  \mathrm{Off}(i,j) &= c - i - S(i,j)
  &&\text{(電源オフ)} \label{eq:Off}
\end{align}
したがって $B+I+S+\mathrm{Off}=c$ が恒等的に成立する。

\paragraph{モデル上の含意}
\begin{itemize}
  \item フルバッチ方式: 待機中の完全バッチ数は $\lfloor j/b\rfloor$ であり,
        $b$ 個揃わないサービスは開始されない。
  \item 需要 $\lfloor j/b\rfloor$ に対してアクティブ台数 $i$ が不足する分を
        \emph{即座に} セットアップ開始する (反応的セットアップ)。
  \item 逆に $j$ が減って需要が下がると, セットアップ中のサーバーは
        \emph{即座に} 取り消される ($S$ が関数であるため)。
\end{itemize}

\subsection{生成行列 $Q$}

\code{build\_generator} は全状態を \eqref{eq:base-index} の順に走査し,
以下の遷移を COO 形式で積み上げてから CSR に変換する
(\code{sum\_duplicates} により重複成分は集約)。
遷移を 1 本追加するごとに対角成分に $-(\text{率})$ を加算する実装である
(\code{add\_transition})。率 $\le 0$ および $\mathrm{src}=\mathrm{dst}$ は追加しない。

\begin{enumerate}
\item \textbf{到着} ($j<K$):
  \begin{equation}
    (i,j,F)\ \longrightarrow\ (i,j+1,F'),\qquad \text{率}\ C_1[F,F']
    \quad (\forall F' \text{ s.t. } C_1[F,F']>0).
  \end{equation}
\item \textbf{到着ブロック} ($j=K$): 位相遷移を伴う到着のみ残す。
  \begin{equation}
    (i,K,F)\ \longrightarrow\ (i,K,F'),\qquad \text{率}\ C_1[F,F']
    \quad (F'\ne F).
  \end{equation}
  $F'=F$ 成分は自己ループとなり $Q$ には現れない。
  本リポジトリの実験では $C_1$ は対角行列であるため, $j=K$ での到着は
  $Q$ に一切寄与しない (損失呼; ジョブは棄却され系の状態は変わらない)。
\item \textbf{バッチサービス完了} ($B(i,j)>0$ かつ $j\ge b$):
  \begin{equation}
    (i,j,F)\ \longrightarrow\ (i,j-b,F),\qquad \text{率}\ B(i,j)\,\mu .
  \end{equation}
  すなわち Busy サーバー 1 台あたり完了率 $\mu$, 1 回の完了で $b$ 個同時退去。
\item \textbf{セットアップ完了} ($S(i,j)>0$ かつ $i<c$):
  \begin{equation}
    (i,j,F)\ \longrightarrow\ (i+1,j,F),\qquad \text{率}\ S(i,j)\,\alpha .
  \end{equation}
  セットアップ中 1 台あたり完了率 $\alpha$ の競合である
  (条件 $i<c$ は $S>0$ から従うので冗長なガード)。
\item \textbf{Delayoff タイムアウト} ($I(i,j)>0$ かつ $i>0$):
  \begin{equation}
    (i,j,F)\ \longrightarrow\ (i-1,j,F),\qquad \text{率}\ I(i,j)\,\beta .
  \end{equation}
  アイドル 1 台あたり率 $\beta$ の独立タイマーで電源オフする
  (条件 $i>0$ は $I>0$ から従うので冗長)。
\item \textbf{MMPP 内部位相遷移}:
  \begin{equation}
    (i,j,F)\ \longrightarrow\ (i,j,F'),\qquad \text{率}\ C_0[F,F']\quad (F'\ne F).
  \end{equation}
\end{enumerate}

\paragraph{帯構造}
到着は $+D$, バッチ完了は $-bD$ のインデックス変位を与えるので,
下帯幅 $p=bD$, 上帯幅 $q=D$ ($C_1$ が対角の場合。非対角成分があれば
$q=D+D_M-1$)。$i$ の増減は変位 $\pm D_M$, 位相遷移は $\pm 1$ 以内で,
いずれも帯の内側である。標準設定 $(c,K,b,D_M)=(20,200,5,2)$ では
$D=42$, $N=8442$, $p=210$, $q=42$ となる。

% =====================================================================
\section{Predictive モデル}
\label{sec:pred}
% =====================================================================

\subsection{拡張の動機と 2 つの機構}

Predictive モデルは \textbf{制御器が MMPP 位相 $F$ を観測できる}
と仮定し (バースト到来の完全予知), 以下 2 機構を追加する。

\begin{description}
\item[(P1) 事前セットアップ] 位相が $0\to1$ に遷移した瞬間
  (通常位相からバースト位相への切替) に, 稼働台数 ($i+s$) が
  $n_{\mathrm{target}}$ に届くよう不足分のサーバーを一斉にセットアップ開始する。
\item[(P2) Delayoff 加速] 通常位相 $F=0$ の間はアイドルサーバーの
  Delayoff 率を $\beta\to\gamma\beta$ に加速し, 早めに電源を落とす。
  バースト位相 $F=1$ では $\beta$ のまま。
\end{description}

実装上, 位相番号 $0$ (通常) と $1$ (バースト) は\textbf{ハードコード}されている
(\code{generator.py} の \code{if F == 0 and Fp == 1}, \code{beta\_F = gamma*beta if F == 0}).
したがって $D_M=2$ の対称 2 位相 MMPP を前提とした実装である。

\subsection{状態空間}

事前セットアップにより $s$ はキューの実需要 $S(i,j)$ を超えうるので,
$s$ を明示的な状態次元に昇格させる。
\begin{equation}
  \mathcal{S}^{P}=\{(i,s,j,F): 0\le i\le c,\ 0\le s\le c-i,\ 0\le j\le K,\ 0\le F<D_M\}.
\end{equation}
拘束は $i+s\le c$ のみ。$(i,s)$ の組を $i$ 昇順・同一 $i$ 内で $s$ 昇順に
列挙して線形番号 $\iota(i,s)$ を与え, $N_{IS}=(c+1)(c+2)/2$ として
\begin{equation}
  \mathrm{idx}(i,s,j,F)=\bigl(j\,N_{IS}+\iota(i,s)\bigr)D_M+F,\qquad
  N^{P}=N_{IS}(K+1)D_M .
  \label{eq:pred-index}
\end{equation}
これはベースの \eqref{eq:base-index} の $i$ の位置に $\iota(i,s)$ を入れた
$j$-major の順序で, レベル $j$ の状態が連続したブロック
$[\,jN_{IS}D_M,\ (j+1)N_{IS}D_M\,)$ になる。
番号付け・逆変換・レベル単位の取り出しは \code{mmpp\_predictive/state\_space.py}
(\code{state\_to\_idx}, \code{idx\_to\_state}, \code{level\_slice}, \code{pi\_by\_level})
に集約している。
標準設定では $N^P = 231\times201\times2 = 92{,}862$。
Busy/Idle は \eqref{eq:B}, \eqref{eq:I} と同じ $(i,j)$ の関数を用いる。

\subsection{事前セットアップ台数}

\begin{equation}
  \Delta = \max(n_{\mathrm{target}}-i-s,\ 0),\qquad
  \Delta_{\eff} = \min(\Delta,\ c-i-s).
  \label{eq:delta}
\end{equation}
第 2 式はオフサーバー不足によるクリップである
(\code{setup\_target\_delta})。$i+s\le c$ の下で $\Delta_{\eff}\ge0$。

\subsection{生成行列 $Q^{P}$}

\code{mmpp\_predictive/generator.py} は \eqref{eq:pred-index} の全名目状態を
走査し, 以下を積み上げる。$s_{\mathrm{need}}(i,j):=S(i,j)$
(ベースモデルの \eqref{eq:S} をそのまま再利用) と書く。

\begin{enumerate}
\item \textbf{到着} ($j<K$): $j'=j+1$ とし
  \begin{equation}
    s' = \begin{cases}
      s+1 & \text{if } s_{\mathrm{need}}(i,j')>s \text{ かつ } s<c-i,\\
      s   & \text{otherwise,}
    \end{cases}
    \qquad
    (i,s,j,F)\to(i,s',j',F'),\ \text{率}\ C_1[F,F'].
  \end{equation}
  すなわち反応的セットアップは $j$ の増加に追従して\textbf{最大 1 台だけ}増える。
\item \textbf{到着ブロック} ($j=K$): $(i,s,K,F)\to(i,s,K,F')$, 率 $C_1[F,F']$, $F'\ne F$。
\item \textbf{セットアップ完了} ($s\ge1$):
  \begin{equation}
    (i,s,j,F)\to(i+1,s-1,j,F),\qquad \text{率}\ s\,\alpha .
  \end{equation}
  $i+s\le c,\ s>0$ より $i<c$ は自動的に満たされる。
\item \textbf{バッチサービス完了} ($B(i,j)>0$ かつ $j\ge b$): $j''=j-b$ とし
  \begin{equation}
    s''=\begin{cases} s-1 & \text{if } s_{\mathrm{need}}(i,j'')<s,\\ s & \text{otherwise,}\end{cases}
    \qquad (i,s,j,F)\to(i,s'',j'',F),\ \text{率}\ B(i,j)\mu .
  \end{equation}
  すなわち需要が下がると\textbf{最大 1 台だけ}セットアップを取り消す。
  これは事前セットアップで積んだ分にも適用される
  (事前セットアップは退去イベントごとに 1 台ずつ解消されていく)。
  ただし選択肢 \code{protect\_presetup} が真のときは \S\ref{sec:protect} の規則に置き換わる。
\item \textbf{MMPP 位相遷移} ($F'\ne F$, 率 $C_0[F,F']$):
  \begin{equation}
    (i,s,j,F)\to
    \begin{cases}
      (i,\ s+\Delta_{\eff},\ j,\ 1) & (F=0,\ F'=1) \quad\text{← 機構 (P1)}\\
      (i,\ s,\ j,\ F')              & \text{otherwise}
    \end{cases}
  \end{equation}
\item \textbf{Delayoff} ($I(i,j)>0$ かつ $i\ge1$):
  \begin{equation}
    (i,s,j,F)\to(i-1,s,j,F),\qquad
    \text{率}\ I(i,j)\,\beta_F,\quad
    \beta_F=\begin{cases}\gamma\beta & F=0\\ \beta & F\ne0\end{cases}
    \quad\text{← 機構 (P2)}
  \end{equation}
\end{enumerate}

\subsection{ベースモデルへの帰着}

$n_{\mathrm{target}}=0$ かつ $\gamma=1$ のとき, \eqref{eq:delta} より
$\Delta_{\eff}\equiv0$ であり, また到着/サービス完了での $s$ の $\pm1$ 追従は
$s=S(i,j)$ を帰納的に保存する。したがって Predictive モデルは
ベースモデルと厳密に一致する。これは
\code{tests/test\_predictive\_baseline\_consistency.py} で
4 指標について許容誤差 $10^{-8}$ で検証されている
(DES 側も \code{tests/test\_predictive\_sim\_baseline\_consistency.py})。

なお $n_{\mathrm{target}}=0$ だけでは $\gamma$ が効くため一致しない。
実験 5 で $n_{\mathrm{target}}=0$ の点を取っているが, そこは
$\gamma=5.0$ が有効な状態である。

\subsection{選択肢: 事前セットアップの保護 (\code{protect\_presetup})}
\label{sec:protect}

上の規則 4 では, バッチ完了のたびに需要 $S(i,j-b)$ を超えたセットアップが
1 台取り消される。これは (P1) で起動を始めた分にも適用されるため,
バースト中に $i+s$ が $n_{\mathrm{target}}$ を下回りうる。
これは (P1) の意図 (稼働中+起動中を $n_{\mathrm{target}}$ にそろえる) と合わないので,
\code{PredictiveModelParameters} に真偽値の選択肢 \code{protect\_presetup} を設けた。
\textbf{既定値は偽}で, そのときは規則 4 のまま (既存の結果は変わらない)。
真のときはバッチ完了の取り消し判定を次のように変える。
\begin{equation}
  \mathrm{cancel}(i,s,j,F)=
  \begin{cases}
    \text{偽} & \text{if } F=1 \text{ かつ } i+s\le n_{\mathrm{target}},\\
    \bigl[\,s_{\mathrm{need}}(i,j-b)<s\,\bigr] & \text{otherwise,}
  \end{cases}
  \qquad s''=s-\mathbf{1}[\mathrm{cancel}(i,s,j,F)].
  \label{eq:protect}
\end{equation}
すなわちバースト位相の間は, 稼働中+起動中が $n_{\mathrm{target}}$ 以下である限り
セットアップを取り消さない。通常位相 $F=0$ では規則 4 と同じである。
判定は \code{mmpp\_predictive/state\_space.py} の \code{setup\_cancelled} に
まとめ, 生成行列と指標 (\code{setup\_cancel\_rate}) の両方がこれを使う。
DES (\code{mmpp\_predictive\_sim/simulator.py}) には同じ規則を独立に実装した。

\paragraph{ベースモデルへの帰着が保たれる理由}
$n_{\mathrm{target}}=0$ のとき, 保護の条件 $i+s\le 0$ は $(i,s)=(0,0)$ でしか
成り立たない。$s=0$ では $s_{\mathrm{need}}\ge0=s$ なので取り消しはもともと起こらず,
\eqref{eq:protect} は規則 4 と同じ遷移を与える。したがって
$n_{\mathrm{target}}=0$ では \code{protect\_presetup} の値によらず
$Q^P$ は同一であり, $\gamma=1$ と合わせてベースモデルに厳密に一致する。
\code{tests/test\_protect\_presetup.py} で, $n_{\mathrm{target}}=0$ での $Q^P$ の
完全一致と, $\gamma=1$ でのベースモデルとの指標の一致 (相対誤差 $10^{-10}$) を確認している。

% =====================================================================
\section{性能指標}
\label{sec:metrics}
% =====================================================================

定常分布 $\pi$ から次を計算する
(\code{mmpp/metrics.py}, \code{mmpp\_predictive/metrics.py})。
以下 $\sum_{\text{state}}$ はベースでは $(i,F)$, Predictive では $(i,s,F)$ にわたる和である。

\begin{align}
  P_\blk &= \Pr(j=K) = \sum_{\text{state}} \pi(\cdot,K,\cdot)
  &&\text{(時間平均ブロック確率)} \\
  \lambda_\eff &= \sum_{j<K}\sum_{\text{state}} \pi(\cdot,j,F)\,\lambda_F
  &&\text{(実効到着率)} \label{eq:lambda_eff}\\
  P_\blk^{\mathrm{arr}} &= 1-\frac{\lambda_\eff}{\bar\lambda}
  &&\text{(到着平均ブロック確率, 定義)} \label{eq:parr}\\
  &= \frac{1}{\bar\lambda}\sum_{\text{state}} \pi(\cdot,K,F)\,\lambda_F
  &&\text{(実装で用いる等価式)} \label{eq:parr_stable}\\
  E[j] &= \sum_{j} j \sum_{\text{state}} \pi(\cdot,j,\cdot) \\
  E[B] &= \sum \pi\, B(i,j),\quad
  E[I] = \sum \pi\, I(i,j) \\
  E[S] &= \begin{cases}
    \sum \pi\, S(i,j) & \text{(ベース: 関数から導出)}\\
    \sum \pi\, s      & \text{(Predictive: 状態変数 $s$ の期待値)}
  \end{cases}\\
  E[\mathrm{Off}] &= c-E[B]-E[I]-E[S] \\
  E[W] &= E[j]/\lambda_\eff \qquad\text{(Little の公式)} \\
  \rho_{\mathrm{server}} &= E[B]/c
\end{align}

\paragraph{$P_\blk \ne P_\blk^{\mathrm{arr}}$ について}
MMPP では位相分布が $j$ に依存するため, 時間平均と到着平均は一致しない。
Poisson 到着 ($D_M=1$) では PASTA により一致する。
\eqref{eq:parr} と \eqref{eq:parr_stable} の数学的等価性は,
$\pi$ の位相マージナルが MMPP 位相定常分布に一致する事実
$\sum_{i,j}\pi(i,j,F)=\varpi_F$ から
\begin{equation}
  \bar\lambda-\lambda_\eff
  =\sum_F \lambda_F\Bigl[\sum_{i,j}\pi(i,j,F)-\sum_{j<K,i}\pi(i,j,F)\Bigr]
  =\sum_{i,F}\pi(i,K,F)\lambda_F
\end{equation}
として導かれる。定義式 \eqref{eq:parr} は $1-(\text{1 に極めて近い値})$ の
引き算を含むため $P_\blk^{\mathrm{arr}}\lesssim10^{-13}$ で桁落ちし負値も出る。
そこで\textbf{実装では引き算を含まない \eqref{eq:parr_stable} だけで計算する}
(\code{arrival\_blocking\_probability\_stable})。
$\pi$ の相対精度 (GTH 法) がそのまま指標に伝わり, $10^{-77}$ 程度の値でも正値かつ
$\rho$ に単調な値が得られる。\eqref{eq:parr} を直接計算する旧メソッド
\code{arrival\_blocking\_probability} はライブラリから削除した。
CSV の出力列は \code{P\_block\_arrival\_stable} のみで,
図・改善率・\code{evaluate\_predictive\_comparison.py} の集計もこの列を用いる。
既存の CSV には引き算で計算した旧列 \code{P\_block\_arrival} が残っているが,
これを読むのは桁落ちを示す検証用スクリプト
(\code{compare\_gth\_vs\_splu\_results.py}, \code{make\_gth\_vs\_splu\_figures.py},
\code{check\_results\_integrity.py}) だけである。
\code{validate\_gth\_v6\_base.py} は比較のために \eqref{eq:parr} をスクリプト内で計算する。

\paragraph{電力コストと ERP}
2 種類のコスト関数が実装されている。
\begin{align}
  \mathrm{Cost}_{\mathrm{generic}} &=
  1.0\,E[B]+0.6\,E[I]+1.0\,E[S]+0.0\,E[\mathrm{Off}]
  \qquad(\code{energy\_cost}) \\
  \mathrm{Cost} &= C_b E[B]+C_s E[S]+C_i E[I],\qquad
  \begin{aligned}
    C_b &= 0.04419\,b+0.15503,\\ C_s &= C_b,\\ C_i &= 0.6\,C_b
  \end{aligned}
  \qquad(\code{energy\_cost\_paper})\\
  \mathrm{ERP} &= E[W]\times \mathrm{Cost}
  \qquad(\code{erp\_paper})
\end{align}
後者は先行研究 (Le-Anh \& Phung-Duc 2025) 5.2 節の式で, オフサーバーの
電力消費を $0$ とする。\textbf{全数値実験は \code{energy\_cost\_paper} /
\code{erp\_paper} を用いている} ($b=5$ では $C_b=C_s=0.37598$, $C_i=0.225588$)。
Predictive モデルでは $C_b,C_s,C_i$ をパラメータとして上書き可能
(既定値は上式と同一)。

\paragraph{事前セットアップ (P1) の診断指標}
Predictive モデルでは次も計算する (\code{all\_metrics} のキー名)。
位相 $0\to1$ の率 $\sigma_{01}$ は $C_0[0,1]$ から取る。
\begin{align}
  \code{p1\_fire\_rate} &= \sigma_{01}\sum_{i,s,j}\pi(i,s,j,0)\,\mathbf{1}[\Delta_\eff(i,s)>0]
  &&\text{(P1 の発動率)}\\
  \code{p1\_launch\_rate} &= \sigma_{01}\sum_{i,s,j}\pi(i,s,j,0)\,\Delta_\eff(i,s)
  &&\text{(P1 で起動を始める台数の率)}\\
  \code{setup\_completion\_rate} &= \sum \pi(i,s,j,F)\,s\,\alpha = \alpha E[S]\\
  \code{setup\_cancel\_rate} &= \sum \pi(i,s,j,F)\,B(i,j)\,\mu\,
     \mathbf{1}[\mathrm{cancel}(i,s,j,F)]
  &&\text{(取り消しの判定は \eqref{eq:protect})}\\
  \code{rho\_B} &= \lambda_1/(c\,b\,\mu)
  &&\text{(バースト位相の負荷; \code{build\_mmpp} では } \rho(1+\delta))
\end{align}
セットアップの開始 (到着による反応的開始と P1) と終了 (完了と取り消し) の率は
定常状態で釣り合う。この恒等式もテストで確認している。

\paragraph{検証恒等式}
実験スクリプトは各点で次の 2 つを診断出力する。
\begin{align}
  E[B]+E[I]+E[S]+E[\mathrm{Off}] &= c &&\text{(保存則)}\\
  \lambda_\eff &= b\,\mu\,E[B] &&\text{(流量バランス)}
\end{align}

% =====================================================================
\section{定常分布の数値解法}
\label{sec:solver}
% =====================================================================

$\pi Q=\mathbf{0}$, $\pi\mathbf{e}=1$ を解く。既定ソルバーは \textbf{GTH 法}。

\subsection{GTH 法 (既定)}

\code{mmpp/gth\_solver.py}。Grassmann--Taksar--Heyman 法
(O'Cinneide 1993, 1996 の GTH1--GTH3) を帯格納で実装する。
生成行列 $G$ の\textbf{非対角成分のみ}を使い, \textbf{引き算を一切行わない}。

\begin{align}
  &\text{GTH1:} && \alpha_k=\sum_{j>k} g_{kj} \\
  &\text{GTH2:} && g_{ij}\leftarrow g_{ij}+\frac{g_{ik}g_{kj}}{\alpha_k}
  \qquad (i,j>k,\ i\ne j) \\
  &\text{GTH3:} && a_{N-1}=1,\quad
  a_k=\frac{1}{\alpha_k}\sum_{i>k} a_i g_{ik},\quad
  \pi=\frac{a}{\sum_k a_k}
\end{align}
全演算が非負値の加算・乗算のみなので桁落ちが原理的に発生せず,
$\pi$ の各成分が\textbf{相対精度}で求まる。

\paragraph{帯格納と計算量}
$\mathrm{Bnd}[r,\,c-r+p]=g_{rc}$ (非対角のみ) に展開する。
帯幅 $(p,q)$ は行列の非零パターンから実測する (\code{bandwidths})。
記憶量 $O(N(p+q))$, 演算量 $O(Npq)$。前進消去 (GTH1+GTH2) と
後退代入 (GTH3) を numba \code{@njit(cache=True)} で高速化し,
標準設定 $(c,K,b)=(20,200,5)$ で約 $0.17$ 秒/点。
ピボット交換は行わず, 更新は帯の内側に限定されるため帯外フィルインは生じない。

\paragraph{前提条件チェック}
GTH は「全状態が最終状態 $N-1$ へ到達可能」であることを要する
(既約ならば成立)。\code{check\_reachability} が $Q$ の逆向きグラフ上で
$N-1$ から BFS して検証し, 満たさなければ \code{ValueError} を送出する。

\paragraph{LU 分解}
\code{return\_factors=True} で掃引後行列 $S$ と $\alpha$ を返し,
\code{build\_LU} が $L$ (単位下三角, $L_{ik}=-S[i,k]/\alpha_k$) と
$U$ ($U_{kk}=\alpha_k$, $U_{N-1,N-1}=0$, $U_{kj}=-S[k,j]$) を構成する。
$-G=LU$ が成り立つ (検証 V4)。

\subsection{Predictive モデルでの到達可能部分空間への縮約}

\eqref{eq:pred-index} の名目状態空間は $i+s\le c$ の全組合せを含むが,
実際に到達可能な状態はその真部分集合である (特に $n_{\mathrm{target}}=0$ では
$s=S(i,j)$ に固定される)。到達不能状態を残すと $Q$ が可約になり
LU も GTH も破綻する。そこで \code{mmpp\_predictive/solver.py} は
\begin{enumerate}
  \item 対角成分を除いた隣接構造上で, 起点
        $\mathrm{idx}=0 \Leftrightarrow (i,s,j,F)=(0,0,0,0)$ から
        \code{breadth\_first\_order} で到達集合を求める,
  \item 到達集合を元のインデックス順にソートして $Q$ を部分行列に縮約,
  \item ベースの \code{solve\_stationary} をそのまま適用,
  \item 得た解を名目サイズのベクトルに散布 (到達不能状態には確率 $0$)
\end{enumerate}
という手順をとる。縮約後の帯幅 $(p,q)$ を実測し,
$N p q > 5\times10^8$ なら \code{RuntimeWarning} を出す。

\textbf{帯幅についての注意}: 索引 \eqref{eq:pred-index} は $j$ が最上位なので,
セットアップ完了 $(i,s)\to(i+1,s-1)$ と Delayoff $(i,s)\to(i-1,s)$ は
同じレベル $j$ の中に収まり, レベルをまたぐ遷移は到着 ($j\to j+1$) と
バッチ完了 ($j\to j-b$) だけになる。したがって帯幅は $K$ によらない。
標準設定 ($c=20$, $b=5$, $n_{\mathrm{target}}=10$, $\gamma=5$, $\rho=0.7$) での
縮約後の実測値は次のとおり。
\begin{center}
\begin{tabular}{rrrrr}
\toprule
$K$ & 縮約後 $N$ & $p$ & $q$ & $Npq$ \\
\midrule
100  & 8{,}092  & 760 & 152 & $9.3\times10^{8}$ \\
200  & 12{,}292 & 760 & 152 & $1.4\times10^{9}$ \\
400  & 20{,}692 & 760 & 152 & $2.4\times10^{9}$ \\
1000 & 45{,}892 & 760 & 152 & $5.3\times10^{9}$ \\
\bottomrule
\end{tabular}
\end{center}
計算量 $O(Npq)$ は $N$ に線形 ($K$ に線形) である。
ただし 1 レベルあたりの状態数がベースより多いため $p,q$ はベース
($p=210$, $q=42$) より大きく, 標準設定でも $Npq>5\times10^8$ となって警告は出る。
旧索引 ($\iota(i,s)$ 最上位, $j$ 中位) ではセットアップ完了と Delayoff が
$(K+1)D_M$ 単位のブロックを跳び越すため帯幅が $K$ に比例して増え
($K=200$ で $(p,q)=(1042,942)$, $Npq\simeq1.2\times10^{10}$;
$K=1000$ で $(2642,2542)$, $Npq\simeq3.1\times10^{11}$),
計算量は $K^3$ に比例していた。
索引の変更による GTH の計算時間 (\code{solve\_stationary} 全体, 2 コア環境) と
帯格納メモリは, $K=200$ で $3.3$\,s $\to$ $0.86$\,s, $186$\,MB $\to$ $86$\,MB,
$K=1000$ で $36$\,s $\to$ $2.3$\,s, $1.8$\,GB $\to$ $320$\,MB である。

\subsection{疎 LU 分解 (\code{solver="splu"}, 旧既定)}

$Q\mathbf{e}=\mathbf{0}$ より $\mathrm{rank}(Q)=N-1$。
$x_N=1$ と固定し, $A=(Q^{\mathsf T})[0{:}N-1,\,0{:}N-1]$,
$u=(Q^{\mathsf T})[0{:}N-1,\,N-1]$ として $Ay=-u$ を
\code{scipy.sparse.linalg.spsolve} で解き, $x=(y,1)$ を正規化する。
$-A$ は非特異 $M$ 行列なので非負一意解をもつ。
ただし引き算を含む通常の LU であり絶対誤差が $\sim10^{-17}$ で頭打ちになるため,
$\pi$ の成分が $10^{-15}$ を下回る領域では相対誤差が破綻する (負値も出る)。
反復法 (GMRES 等) は本モデルの剛性 ($\kappa\sim10^3$) と
動的レンジ ($\kappa_\pi\sim10^{10}$--$10^{20}$) では不適と判断して採用していない。

\subsection{解の事後検証}

\code{solve\_stationary} は解を返す前に
$\|\pi Q\|_\infty\le \mathrm{tol}$, $|\sum\pi-1|\le\mathrm{tol}$,
$\min\pi\ge-\mathrm{tol}$ (既定 $\mathrm{tol}=10^{-8}$) を検査し,
違反時は \code{RuntimeError}。通過後, 数値誤差による微小負値を $0$ に
クリップして再正規化する。

\subsection{検証結果の要旨}

\code{results\_gth/VALIDATION\_REPORT.md} に詳細がある。要点のみ:
\begin{itemize}
  \item M/M/1/K 解析解および O'Cinneide (1996) \S4 の 20 状態例で
        相対誤差 $0.0$ (ビット一致)。同例を通常の Gauss 消去で解くと
        相対誤差 $5.71$, 最小成分 $-2.9\times10^{-17}$ (負の確率)。
  \item mpmath 50 桁の真値との比較 (ベース): $P_\blk$ が $10^{-65}$ でも
        GTH の相対誤差 $\sim10^{-15}$。splu は $P_\blk\lesssim10^{-12}$ で破綻
        (最悪ケースで真値と 29 桁ずれ)。
  \item Predictive モデル (小規模設定 $c=5,K=60,b=2,n_{\mathrm{target}}=3,\gamma=5$,
        縮約後 $N=788$) でも GTH 相対誤差最大 $5.78\times10^{-15}$。
  \item 既存 30 実験ファイルで GTH vs splu を突き合わせ, ブロック確率以外は
        相対差 $\lesssim2\times10^{-12}$。相対差が $10^{-4}$ を超えた 14 点は
        \emph{すべて}定義式 \eqref{eq:parr} の引き算で計算した旧列
        \code{P\_block\_arrival} であり,
        時間平均 $P_\blk$ では 1 点も超過しない。これが原因が指標式の
        桁落ちであることの証拠である。
\end{itemize}

% =====================================================================
\section{離散事象シミュレーション (DES)}
\label{sec:des}
% =====================================================================

理論実装の独立検証用。\code{mmpp\_sim/}, \code{mmpp\_predictive\_sim/}。

\subsection{方式}

競合指数分布法 (Gillespie 法)。現状態から発生しうる全遷移を列挙し,
\begin{enumerate}
  \item 合計率 $R=\sum_k r_k$ で滞在時間 $\Delta t\sim\mathrm{Exp}(R)$ を生成,
  \item $u\sim\mathcal{U}(0,R)$ で率に比例する確率でイベント種別を選択
\end{enumerate}
する。全事象が指数分布 (無記憶) なので真の CTMC 軌道と同分布である。

イベント種別は 5 種: \code{ARRIVAL}, \code{SERVICE\_COMPLETION},
\code{SETUP\_COMPLETION}, \code{DELAYOFF\_TIMEOUT}, \code{PHASE\_TRANSITION}。
Predictive 版では \code{PHASE\_TRANSITION} が $F:0\to1$ のとき
同一イベント内で $s\mathrel{+}=\Delta_{\eff}$ を適用し, その発動回数と
台数を別途集計する。

\paragraph{$j=K$ での到着の扱い (重要)}
DES のイベントメニューは $j=K$ でも率 $C_1[F,F']$ の到着を
\emph{すべて} (自己ループ $F'=F$ を含めて) 候補に入れる。
これは $P_\blk^{\mathrm{arr}}=\text{ブロック数}/\text{到着試行数}$ の
統計を正しく取るためである。理論側の $Q$ では自己ループが消えているため
一見不整合だが, 自己ループを選んだ場合は状態が変わらず,
その滞在時間 $\Delta t$ も \code{record\_state} で加算されるので,
無記憶性より連続する自己ループ分の合計滞在時間は真の
$\mathrm{Exp}(R_{\text{実遷移}})$ と同分布になる。
したがって時間重み付き統計は不偏である。

\subsection{統計収集}

\code{SimStats} が時間重み付き統計を蓄積する: 状態別累積滞在時間,
総記録時間 $T$, 到着試行数, ブロック数, 退去ジョブの滞在時間総和と総数。

\paragraph{滞在時間 (sojourn time) の帰属近似}
\code{SERVICE\_COMPLETION} 時に FIFO キューから最古の $b$ 個の到着時刻を
pop して滞在時間を計算する。物理的には $B(i,j)>1$ のとき完了するバッチは
一様ランダムであり最古とは限らないので, これは厳密ではない。
しかし任意の 1 対 1 対応 $\sigma$ について
$\sum_k(d_k-a_{\sigma(k)})=\sum_k d_k-\sum_k a_k$ であり総和は帰属方法に
依らないため, 集計平均 $E[W]$ は影響を受けず Little の公式との整合も保たれる。
FIFO を使う理由は実装の単純さのみである。

\subsection{実行手順と信頼区間}

\begin{itemize}
  \item 初期状態 $(i,j,F)=(0,0,0)$ (Predictive は $(0,0,0,0)$), 時刻 $0$。
  \item \code{warmup\_events} 回は統計を記録せず実行 (既定 $10^5$)。
        このときの実時間を \code{warmup\_duration} に記録する。
  \item 続く \code{measurement\_events} 回を記録 (既定 $10^6$)。
  \item \textbf{独立レプリケーション法}: \code{run\_replications} が
        $n_{\mathrm{reps}}$ 個 (既定 20) の独立シミュレータを
        シード $\mathrm{seed}_0+k$ で実行する。
  \item 点推定は全レプリケーションの統計を\textbf{合算}した
        \code{SimStats} から計算する。
  \item 95\% 信頼区間はレプリケーションごとの指標値から
        $t$ 分布近似 ($t_{0.975,n-1}\cdot s/\sqrt{n}$) で求める。
        非有限値のレプリケーションは除外し警告を出す。
\end{itemize}

\paragraph{ウォームアップ妥当性診断}
\code{check\_warmup\_duration} は
$\mathrm{warmup\_duration}\times\beta<5$ のとき
「最も遅い時定数 $1/\beta$ に対しウォームアップが不足」と
\code{RuntimeWarning} を出す。

\subsection{DES 側の指標}

\begin{align}
  P_\blk &= \frac{\sum_{\text{states with }j=K}\tau_{\text{state}}}{T},\qquad
  E[j] = \frac{\sum_{\text{state}} j\,\tau_{\text{state}}}{T},\\
  P_\blk^{\mathrm{arr}} &= \frac{\#\text{blocked}}{\#\text{attempts}},\qquad
  \lambda_\eff = \frac{\#\text{attempts}-\#\text{blocked}}{T},\\
  E[W] &= \frac{\sum \text{sojourn}}{\#\text{departures}}.
\end{align}
DES の $P_\blk^{\mathrm{arr}}$ はブロック数 / 到着試行数による経験的推定量であり,
理論側の定義式 \eqref{eq:parr} の引き算とは無関係なのでそのまま残している
(CSV・辞書のキーも \code{P\_block\_arrival})。
理論値と比較する場合は, 理論側の値として \eqref{eq:parr_stable}
(\code{P\_block\_arrival\_stable}) を用いる。
その他 ($E[B],E[I],E[S],E[\mathrm{Off}],\rho,\mathrm{Cost},\mathrm{ERP}$) は
理論側と同一の定義式を時間重み付き平均で評価する。
Predictive 版では $E[S]$ を状態変数 $s$ の時間平均として直接計算する。

% =====================================================================
\section{数値実験の設計}
\label{sec:experiments}
% =====================================================================

\subsection{バースト MMPP の構成}

全実験の到着過程は\textbf{対称 2 位相 MMPP} である
(\code{scripts/\_mmpp\_burst.py} の \code{build\_mmpp})。
トラフィック強度 $\rho$, 振幅 $\delta$, 時定数 $\sigma$ から
\begin{align}
  \bar\lambda &= \rho\, c\, b\, \mu, \\
  \lambda_0 &= \bar\lambda(1-\delta) \quad\text{(通常位相)}, \qquad
  \lambda_1 = \bar\lambda(1+\delta) \quad\text{(バースト位相)}, \\
  C_1 &= \begin{pmatrix}\lambda_0 & 0\\ 0 & \lambda_1\end{pmatrix},\qquad
  C_0 = \begin{pmatrix}-(\sigma+\lambda_0) & \sigma\\
                        \sigma & -(\sigma+\lambda_1)\end{pmatrix}.
\end{align}
位相遷移率は $0\to1$, $1\to0$ ともに $\sigma$ (各位相の平均滞在時間 $1/\sigma$)。
$C_1$ が\textbf{対角}であることから\textbf{到着は位相を変えない}。
定常位相確率は $(1/2,1/2)$ なので, $\delta,\sigma$ を変えても
$\bar\lambda=\rho c b\mu$ は不変である (バースト性のみを独立に変えられる設計)。

$\rho$ は「飽和スループット $cb\mu$ に対する到着率の比」として定義されている
点に注意 ($\rho=\bar\lambda/(cb\mu)$)。これは実測のサーバー利用率
$\rho_{\mathrm{server}}=E[B]/c$ とは別物である。

\paragraph{到着間隔の変動係数 CV}
図のキャプション用に \code{compute\_interarrival\_cv} が
phase-type 分布の 2 次モーメント
\begin{equation}
  E[X^2]=2\,\boldsymbol{p}_{\mathrm{arr}}(-C_0)^{-2}\mathbf{e},\qquad
  \boldsymbol{p}_{\mathrm{arr}}=\frac{\boldsymbol{\varpi}C_1}{\bar\lambda}
\end{equation}
(Fischer \& Meier-Hellstern 1993, "The MMPP cookbook") から
$\mathrm{CV}=\sqrt{E[X^2]-E[X]^2}/E[X]$, $E[X]=1/\bar\lambda$ を計算する。

\subsection{共通パラメータ}

\begin{center}
\begin{tabular}{lll}
\toprule
項目 & 値 & 備考 \\
\midrule
\code{BASELINE} & $c=20,\ K=200,\ b=5,\ \mu=1.0,\ \alpha=0.1,\ \beta=0.005$
  & $1/\alpha=10$ s, $1/\beta=200$ s \\
weak   & $(\delta,\sigma)=(0.3,\ 1.0)$ & \\
medium & $(\delta,\sigma)=(0.6,\ 0.1)$ & 既定のバースト水準 \\
strong & $(\delta,\sigma)=(0.9,\ 0.01)$ & \\
\code{ALPHA\_LEVELS} & $\alpha\in\{0.1,\ 1.0,\ 10.0\}$ & 実験 2,3,5,6 で走査 \\
Predictive 既定値 & $n_{\mathrm{target}}=10,\ \gamma=5.0$ & \\
\code{ALPHA\_GAMMA\_MAP} & $0.1\!\mapsto\!1.0,\ 1.0\!\mapsto\!100.0,\ 10.0\!\mapsto\!1000.0$
  & \code{--alpha-gamma-map} 指定時 \\
\bottomrule
\end{tabular}
\end{center}

パラメータ規模は先行研究 (Le-Anh \& Phung-Duc 2025) Fig.~8/Fig.~10 に合わせている。
$K=200$ は $cb\times4=400$ を計算コストの都合で半分にした値である。

\code{--alpha-gamma-map} を指定した場合, $\alpha$ ごとに異なる $\gamma$ を使う。
対応は厳密一致を優先し, 無い場合は $\log_{10}\alpha$ 距離最小の
キーにフォールバックする (\code{resolve\_gamma})。

\subsection{実験 0 / 0-P: 理論とシミュレーションの整合性検証}

\begin{center}
\begin{tabular}{ll}
\toprule
項目 & 設定 \\
\midrule
スクリプト & \code{experiment\_0\_validation.py} / \code{experiment\_0P\_validation.py} \\
パラメータ & \code{BASELINE}, medium バースト ($\delta=0.6,\sigma=0.1$) \\
走査 & $\rho\in[0.3,\,0.95]$ を等間隔 10 点 \\
Predictive & $n_{\mathrm{target}}=10,\ \gamma=5.0$ \\
DES & $n_{\mathrm{reps}}=20$, warmup $10^5$ events, measurement $10^6$ events, seed $42$ \\
比較指標 & $P_\blk,\ E[N],\ E[W],\ \lambda_\eff,\ \rho_{\mathrm{server}},\ \mathrm{Cost}$ の 6 つ \\
合否 & 各 $\rho$ で「理論値が 95\% CI 内」が 6 指標中 5 つ以上を目安 \\
追加診断 & 保存則と流量バランスの最大誤差 \\
出力 & \code{figures/experiment\_0\_validation.\{png,pdf\}} 他 \\
\bottomrule
\end{tabular}
\end{center}

\textbf{実験 0 でこの一致を確認したうえで, 実験 1 以降は理論解析のみで図を作る}
という方針である (走査点数が多くシミュレーションでは計算時間が過大になること,
論文図として理論曲線の滑らかさを優先することが理由)。

検証結果 (GTH 既定下): 実験 0 はベースモデルで全 10 点・全 6 指標が CI 内。
実験 0-P は低 $\rho$ の 2 点 ($\rho=0.30,0.37$) で $E[W],E[N]$ が
わずかに CI を外れた (例: 理論 $E[W]=1.2233$ に対し CI $[1.2161,1.2227]$, 差 0.05\%)。
両ソルバーで理論値が機械精度で同一であることを確認済みのため,
これは GTH 由来ではなく DES 側のバイアス (ウォームアップ長・測定長) である。

\subsection{実験 1--4 / 1-P--4-P: 制御変数の走査 (理論のみ)}

4 指標 $P_\blk^{\mathrm{arr}}$ (式 \eqref{eq:parr_stable} で計算,
CSV 列 \code{P\_block\_arrival\_stable}), $E[W]$,
$\mathrm{Cost}$, $\mathrm{ERP}$ を $2\times2$ サブプロットで出力する。
$P_\blk^{\mathrm{arr}}$ の $y$ 軸は対数スケールで, 下限は設けない
(実験 4/4-P の低 $\rho$・大 $K$ では $10^{-77}$ 程度まで描かれる)。
値の範囲が広い図 (実験 1/1-P/4/4-P と比較図 1 vs 1-P, 4 vs 4-P) では,
対数軸の自動余白で上端が 1 を超えないよう上端を 1 に固定している。走査変数が対数格子のもの
($\beta$: 実験 2/2-P, $\sigma$: 実験 3-B/3-P-B, $\gamma$: 実験 6) は $x$ 軸も対数スケール。

\begin{longtable}{p{0.10\textwidth}p{0.40\textwidth}p{0.42\textwidth}}
\toprule
実験 & 走査変数 & 固定条件 \\
\midrule
\endhead
1 / 1-P & $\rho\in[0.1,0.95]$ 等間隔 15 点 $\times$ 3 バースト水準
  & \code{BASELINE}。1-P は $n_{\mathrm{target}}=10$, $\gamma\in\{5.0,\,1.0\}$ \\
2 / 2-P & $\beta\in[10^{-2},10^{2}]$ 対数 15 点 $\times$ $\alpha\in\{0.1,1,10\}$
  & $c=20,K=200,b=5,\mu=1,\rho=0.7$。メイン medium, 補助 strong
    (\code{--strong-burst})。2-P は $n_{\mathrm{target}}=10$,
    $\gamma=5.0$ または \code{--alpha-gamma-map} \\
3-A / 3-P-A & $\delta\in[0,0.95]$ 等間隔 15 点 $\times$ $\alpha$ 3 水準
  & \code{BASELINE} ($\beta=0.005$ 固定), $\rho=0.7$, $\sigma=0.1$ 固定 \\
3-B / 3-P-B & $\sigma\in[10^{-3},10^{1}]$ 対数 15 点 $\times$ $\alpha$ 3 水準
  & 同上, $\delta=0.6$ 固定 \\
4 / 4-P & $\rho\in[0.1,0.95]$ 15 点 $\times$ $K\in\{100,200,500,1000\}$
  & $c=20,b=5,\mu=1,\alpha=0.1,\beta=0.005$。3 バースト水準それぞれで実行。
    4-P は $n_{\mathrm{target}}=10$, $\gamma\in\{5.0,\,1.0\}$ \\
\bottomrule
\end{longtable}

$\delta$ の上限を $0.95$ にしているのは $\delta\to1$ で $\lambda_0\to0$ となり
数値的な特異に近づくのを避けるためである ($\delta=0$ は Poisson 到着)。

\subsection{実験 5, 6: Predictive 固有パラメータの感度 (Predictive のみ)}

\begin{center}
\begin{tabular}{p{0.08\textwidth}p{0.42\textwidth}p{0.42\textwidth}}
\toprule
実験 & 走査変数 & 固定条件 \\
\midrule
5 & $n_{\mathrm{target}}\in\{0,1,3,4,6,7,9,10,11,13,14,16,17,19,20\}$ (15 点)
  & \code{BASELINE}, $\rho=0.7$, $\gamma=5.0$,
    $\alpha$ 3 水準 $\times$ 3 バースト水準 \\
6 & $\gamma$: $[1,20]$ の対数 15 点のうち $5.54$ を代表値 $5.0$ に置換
  & \code{BASELINE}, $\rho=0.7$, $n_{\mathrm{target}}=10$,
    $\alpha$ 3 水準 $\times$ 3 バースト水準 \\
\bottomrule
\end{tabular}
\end{center}

実験 6 の $\gamma$ 水準は
$\{1.0,\,1.239,\,1.534,\,1.900,\,2.354,\,2.915,\,3.611,\,4.472,\,5.0,\,6.861,\,
8.498,\,10.525,\,13.037,\,16.147,\,20.0\}$。
境界 $\gamma=1$ (加速無効), $\gamma=20$ (拡張走査との接続点),
代表値 $\gamma=5.0$ を必ず含むように構成している。
\code{compare\_experiment\_6\_extended.py} は初期走査 ($\gamma\le20$) と
拡張走査 ($\gamma\in\{20,50,100,500,1000\}$) の CSV を統合して
$\gamma\in[1,1000]$ の依存性を可視化し, 最適 $\gamma$ の位置を特定する。

\subsection{ベース vs Predictive の比較}

\code{compare\_experiment\_\{1,2,3,4\}\_vs\_\{1P,2P,3P,4P\}.py} が
両者の CSV を読み込み, 同一指標をベース (実線) と Predictive (破線) で
重ね描きする。改善率は
\begin{equation}
  \text{improvement [\%]}=\frac{\text{base}-\text{pred}}{\text{base}}\times100
\end{equation}
(小さいほど良い指標を前提とした符号) として標準出力に要約される。

\code{evaluate\_predictive\_comparison.py} は実験 1--4 と 1-P--4-P を
横断的に突き合わせ, 勝率・平均改善率・条件別詳細・パターン抽出を
Markdown で出力する (\code{evaluation\_report*.md})。
ベース実験と Predictive 実験でグリッド密度が異なるため,
連続変数は\textbf{最近傍マッチング} ($\beta,\sigma$ は $\log$ 距離,
$\rho,\delta$ は線形距離), 離散変数 ($\alpha$, $K$, バースト名) は
\textbf{厳密一致}で対応付ける。
勝敗は相対差 $|\text{base}-\text{pred}|/|\text{base}|<10^{-6}$
(\code{TIE\_REL\_TOL}) なら\textbf{引き分け}とする。
Predictive の機構がほとんど働かない点 (弱バースト, $\gamma=1$ など) では
両者の差が丸め誤差 ($\sim10^{-12}$) しかなく, 引き分けにしないと
勝敗が数値誤差で決まるためである。改善率の 0 除算回避以外の足切りは設けず,
$P_\blk^{\mathrm{arr}}$ は「全点」と「ベース値 $\ge10^{-6}$ の点のみ」の
2 通りの勝敗を併記する。$P_\blk^{\mathrm{arr}}$ の効果量は, 平均改善率が極小値どうしの
比で支配されるため出力せず, $\log_{10}(P_{\mathrm{pred}}/P_{\mathrm{base}})$ の中央値と
四分位範囲で要約する (負なら Predictive の方が小さい)。
さらに勝ち・負け・引き分けの数を $\alpha$ 水準別 ($0.1, 1, 10$) とバースト水準別に分けて出力する
(実験 1-P/4-P は $\alpha=0.1$ 固定, 3-P はバースト水準を持たないので走査ごとに別の行)。
\code{\_gmap} と \code{g1.0} を組み合わせた比較 (\code{evaluation\_report\_optimal\_gamma.md}) では
$\alpha=0.1$ の点がすべて $\gamma=1$ (P2 無効) なので, その行は P1 単独の効果として読める。

\code{ALPHA\_GAMMA\_MAP} ($0.1\!\mapsto\!1,\ 1\!\mapsto\!100,\ 10\!\mapsto\!1000$) の
決め方はコミット履歴・スクリプト・コメントのいずれにも記録がない。
実験 6 (初期 + 拡張走査) で ERP を最小にする $\gamma$ は
$\alpha=0.1$: 1 (weak, medium), 20 (strong);
$\alpha=1$: 46 (weak), 187 (medium, strong);
$\alpha=10$: 1000 (全バースト, 走査上端) であり, 桁としては対応するが一致はしない。

\subsection{出力ファイルの命名規約}

\begin{itemize}
  \item ベース: \code{results/experiment\_1.csv},
        \code{experiment\_2\_\{medium,strong\}.csv},
        \code{experiment\_3\_\{delta,sigma\}.csv},
        \code{experiment\_4\_\{weak,medium,strong\}.csv}
  \item Predictive: \code{experiment\_1P\_nt10\_g5.0.csv} のように
        \code{nt}$\langle n_{\mathrm{target}}\rangle$\code{\_g}$\langle\gamma\rangle$ を接尾, \code{--alpha-gamma-map}
        指定時は \code{\_gmap}
  \item 図は \code{figures/} に PNG (dpi 150) と PDF (dpi 100) の両方。
        図中に題名は付けず, 設定情報はコンソールに「キャプション用情報」として
        出力する方針 (論文キャプション側に記載する前提)。
  \item \code{results/} と \code{results\_gth/} の CSV は\textbf{どちらも GTH 法で計算したもの}で,
        対応する 30 ファイルは全指標・全行が完全に同一である
        (\code{results/experiment\_6\_extended/} の 3 ファイルは \code{results/} にのみある)。
        経緯: \code{results\_gth/} は GTH 既定化直後 (2026-09-25, コミット
        \code{8f3cb6f}, \code{f9594ff}) に全実験を再計算した結果で, その時点の
        \code{results/} は splu で計算した旧結果だった。GTH と splu の突き合わせ
        (\code{compare\_gth\_vs\_splu\_results.py}, \code{VALIDATION\_REPORT.md}) は
        この時点の両者で行った。その後 2026-09-28 のコミット \code{186a895} で
        \code{results/} が GTH の結果で上書きされ, \code{results\_gth/} と同一になった。
        したがって \code{results\_gth/} は現在は\textbf{重複}であり, 数値実験の正本は
        \code{results/} である (ディレクトリは検証の記録として残している)。
        現在の両ディレクトリで \code{compare\_gth\_vs\_splu\_results.py} を実行しても
        差は 0 になる。splu 時代の CSV は \code{git show 186a895\^{}:results/<ファイル名>}
        で取り出せる。
  \item \code{figures\_gth/} は到着平均ブロッキング確率を引き算版の式で描いた旧図であり,
        論文・スライドには \code{figures/} を使う (\code{figures\_gth/README.md})。
\end{itemize}

\subsection{補助スクリプト}

\begin{center}
\begin{tabular}{ll}
\toprule
スクリプト & 用途 \\
\midrule
\code{validate\_gth\_mpmath.py} & mpmath 50 桁の真値との比較 (ベース) \\
\code{validate\_gth\_mpmath\_predictive.py} & 同 (Predictive, 小規模設定) \\
\code{validate\_gth\_v6\_base.py} & 実験格子を直接再現して GTH/splu を同一 $Q$ 上で比較 (415 点) \\
\code{compare\_gth\_vs\_splu\_results.py} & \code{results/} (splu 時代) と \code{results\_gth/} の CSV 突き合わせ (現在は両者同一のため差 0) \\
\code{benchmark\_gth\_vs\_splu.py} & 計算時間・帯格納メモリの計測 (\code{results\_gth/benchmark.md}) \\
\code{make\_gth\_vs\_splu\_figures.py} & 精度比較図の生成 \\
\code{check\_results\_integrity.py} & CSV の \code{P\_block\_arrival\_stable} 列の存在・正値性の検査 \\
\code{replot\_from\_csv.py} & 既存 CSV から実験 1--6 の図だけを描き直す (定常分布は再計算しない) \\
\bottomrule
\end{tabular}
\end{center}

% =====================================================================
\section{テストによる妥当性検証}
% =====================================================================

\code{tests/} の主要テストと, それが担保する性質。

\begin{longtable}{p{0.40\textwidth}p{0.55\textwidth}}
\toprule
テスト & 担保する性質 \\
\midrule
\endhead
\code{test\_model.py} & パラメータ検証 (符号規約, MMPP 整合性, 既約性, 警告) \\
\code{test\_state\_space.py} & \eqref{eq:base-index} の全単射性, \code{block\_slice} \\
\code{test\_generator.py} & 行和 $0$, 非対角 $\ge0$, 各遷移の率と行き先 \\
\code{test\_solver.py} & $\pi Q=0$, 正規化, 非負性, 位相マージナル $=\varpi$ \\
\code{test\_gth\_solver.py} & M/M/1/K 解析解, O'Cinneide 例, 帯格納の厳密保存, $-G=LU$ \\
\code{test\_metrics.py} & 保存則, 流量バランス, 各指標の定義 \\
\code{test\_stable\_blocking.py} & \eqref{eq:parr_stable} の正値性 (\code{results/} の全実験格子 + 同じ $\rho$ 格子の縮小モデル) と, 高精度領域で $\bar\lambda-\lambda_\eff$ の直接和との一致 \\
\code{test\_predictive\_state\_space.py} & \eqref{eq:pred-index} の全単射性, $\Delta_{\eff}$ \\
\code{test\_predictive\_generator.py} & Predictive の各遷移 \\
\code{test\_predictive\_baseline\_consistency.py}
  & $n_{\mathrm{target}}=0,\gamma=1$ で理論がベースと一致 (最重要) \\
\code{test\_predictive\_sim\_baseline\_consistency.py}
  & 同じ一致を DES 側でも確認 \\
\code{test\_gth\_predictive.py} & 縮約 $\to$ GTH 経路の正しさ, GTH と splu の一致 \\
\code{test\_protect\_presetup.py}
  & \code{protect\_presetup} が偽で既存の結果が不変, 真でも $n_{\mathrm{target}}=0,\gamma=1$ で
    ベースと一致, 保護でバースト中の $i+s<n_{\mathrm{target}}$ の確率が減ること,
    P1 診断指標 (セットアップの流量保存, DES の発動・取り消し回数との一致) \\
\code{test\_simulator.py}, \code{test\_predictive\_sim\_simulator.py}
  & イベントメニュー, 状態遷移, 統計収集 \\
\code{test\_experiment\_*.py} & 各実験スクリプトのパラメータ格子・CSV 列・図生成 \\
\code{test\_alpha\_gamma\_map.py} & \code{--alpha-gamma-map} のパースとフォールバック \\
\bottomrule
\end{longtable}

% =====================================================================
\section{認識確認のためのチェックリスト}
\label{sec:checklist}
% =====================================================================

実装上の判断が分かれうる点, あるいは「普通に読むと別の解釈をしそうな点」を
列挙する。各項目が意図どおりかご確認いただきたい。

\paragraph{モデル定義}
\begin{enumerate}
\item \textbf{$i$ の意味}: $i$ は「アクティブ (電源が入っていて仕事ができる) サーバー数」で,
      Busy と Idle (Delayoff 計時中) の\emph{両方}を含む。
      セットアップ中のサーバーは $i$ に含まれない。
\item \textbf{フルバッチ}: $b$ 個揃わないとサービスを開始しない。
      Busy 台数は $\min(i,\lfloor j/b\rfloor)$ であり,
      端数 $j \bmod b$ のジョブはサーバーを占有しない。
\item \textbf{ベースモデルのセットアップは瞬時反応}: $S(i,j)$ は関数なので,
      $j$ が増えれば即セットアップ開始, $j$ が減れば即キャンセルされる。
      セットアップ「開始」「中止」自体に遅延やコストはない。
\item \textbf{Delayoff は 1 台ずつ独立タイマー}: 率 $I(i,j)\beta$ で
      アイドル 1 台が落ちる。Delayoff 中のサーバーは即座にジョブを取れる
      (すなわち「オフ」ではなく「低電力待機」)。
\item \textbf{$j=K$ での到着は完全棄却}: $C_1$ が対角の本実験設定では,
      ブロックされた到着は系の状態を一切変えない (再試行なし, キュー外待ちなし)。
\item \textbf{Predictive は位相を完全観測できる}と仮定している。
      予測の誤り (false positive / false negative) はモデル化していない。
\item \textbf{(P1) は位相 $0\to1$ の瞬間のみ}発動する。バースト中に
      $i+s$ が $n_{\mathrm{target}}$ を下回っても追加発動はしない。
\item \textbf{(P2) は位相 $F=0$ のとき常時}有効。$\gamma\beta$ は
      アイドル 1 台あたりの率に掛かる。
\item \textbf{事前セットアップ分の解消}: 退去イベント (バッチサービス完了) 1 回で
      $s$ は最大 1 台しか減らない。逆に $\Delta_{\eff}$ は一度に複数台増やせる。
      この非対称性は意図的か。
\item \textbf{位相番号のハードコード}: $F=0$ を通常, $F=1$ をバーストとして
      \code{generator.py} に直接書かれている。$D_M>2$ や位相順序を変えた
      MMPP には対応していない。
\item \textbf{$\rho$ の定義}: $\rho=\bar\lambda/(cb\mu)$ であり, 系の実利用率ではない。
      Delayoff/セットアップにより実測 $\rho_{\mathrm{server}}=E[B]/c$ は $\rho$ と一致しない。
\end{enumerate}

\paragraph{指標}
\begin{enumerate}\setcounter{enumi}{11}
\item \textbf{到着平均ブロッキング確率は引き算を含まない式 \eqref{eq:parr_stable} だけで計算する}。
      定義式 $1-\lambda_\eff/\bar\lambda$ (引き算版) はライブラリから削除し,
      CSV にも \code{P\_block\_arrival\_stable} 列だけを出力する。
      図・改善率・比較集計もすべてこの列の値であり,
      極小領域 ($\lesssim10^{-13}$) も相対精度を保ったまま図示されている。
      (旧 CSV に残る \code{P\_block\_arrival} 列は桁落ちの検証用にのみ参照する。)
\item \textbf{図の $y$ ラベルは $P_{\blk}$ だが中身は到着平均}である
      (\code{METRIC\_SPECS} のラベルが $P_{\mathrm{block}}$,
      プロットする値は \code{P\_block\_arrival\_stable})。
      時間平均 $P_\blk$ と混同しないよう, 論文側の記述に注意。
\item \textbf{$E[W]$ は系内滞在時間} (待ち時間 + サービス時間) であり,
      Little の公式 $E[W]=E[j]/\lambda_\eff$ で求めている
      ($E[j]$ が待機中 + 処理中の双方を含むため)。
\item \textbf{コスト関数}: 実験は \code{energy\_cost\_paper}
      ($C_b=C_s=0.04419b+0.15503$, $C_i=0.6C_b$, Off は $0$) を使用。
      汎用版 \code{energy\_cost} ($1.0/0.6/1.0/0.0$) は実験では使っていない。
\item \textbf{ERP} は $E[W]\times\mathrm{Cost}$ であり, 小さいほど良い。
\end{enumerate}

\paragraph{数値解析}
\begin{enumerate}\setcounter{enumi}{16}
\item \textbf{既定ソルバーは GTH}。\code{solver="splu"} で旧実装に戻せるが,
      極小ブロック確率領域では splu の値は信用できない (最悪 29 桁の誤り)。
\item \textbf{GTH は対角成分を一切参照しない}。したがって
      \code{build\_generator} が積む対角成分は splu 経路と事後検証
      ($\|\pi Q\|_\infty$) のためだけに存在する。
\item \textbf{Predictive は到達可能部分空間に縮約してから}解く。
      縮約の起点は $\mathrm{idx}=0\Leftrightarrow(i,s,j,F)=(0,0,0,0)$。
      名目状態空間のままでは $Q$ が可約で解けない。
\item \textbf{Predictive の計算コスト}: 索引順序 \eqref{eq:pred-index} は
      $j$ 最上位なので縮約後の帯幅は $K$ によらず, 標準設定で
      $(p,q)=(760,152)$, $Npq\simeq1.4\times10^9$ ($K=200$)。
      計算量は $K$ に線形だが, $p,q$ がベース ($210, 42$) より大きいため
      ベースモデル ($0.17$ 秒/点) よりは重く, $Npq>5\times10^8$ の警告は出る。
\item \textbf{事後検証の許容誤差は $10^{-8}$} (絶対値基準)。
      通過後に微小負値をクリップして再正規化している。
\end{enumerate}

\paragraph{シミュレーション}
\begin{enumerate}\setcounter{enumi}{21}
\item \textbf{DES の到着は $j=K$ でも自己ループを含めて生成する}
      (ブロック率の分母を正しく数えるため)。理論側の $Q$ では自己ループが
      消えているが, 時間重み付き統計は無記憶性により一致する。
\item \textbf{滞在時間は FIFO 帰属の近似}。個別ジョブの滞在時間分布は
      厳密ではないが, 平均 $E[W]$ は厳密 (全単射不変性)。
      滞在時間の\emph{分布}や分位点を論じる場合はこの実装では不適。
\item \textbf{終了条件はイベント数}であり実時間ではない
      (warmup $10^5$, measurement $10^6$ イベントが既定)。
\item \textbf{信頼区間は独立レプリケーション法} ($n=20$) +
      $t$ 分布近似。点推定は全レプリケーション合算値であり,
      レプリケーション平均ではない (バッチサイズが揃っていれば近いが同一ではない)。
\item \textbf{ウォームアップの十分性}は $\mathrm{warmup\_duration}\times\beta\ge5$
      を目安に自動診断している。$\beta=0.005$ では $1/\beta=200$ なので
      実時間 $1000$ 以上が必要。
\end{enumerate}

\paragraph{実験手順}
\begin{enumerate}\setcounter{enumi}{26}
\item \textbf{実験 1 以降はシミュレーションを行っていない} (理論のみ)。
      理論の正当性は実験 0 / 0-P で担保するという構成。
\item \textbf{$\delta,\sigma$ を変えても $\bar\lambda$ は不変}なので,
      実験 3 はバースト性のみを分離して見ている。
\item \textbf{実験 5 の $n_{\mathrm{target}}=0$ 点はベースと一致しない}
      ($\gamma=5.0$ が有効なまま)。ベースとの厳密一致は
      $n_{\mathrm{target}}=0$ かつ $\gamma=1$ のときのみ。
\item \textbf{ベースと Predictive のグリッド密度が異なる}ため,
      \code{evaluate\_predictive\_comparison.py} は最近傍マッチングで
      対応付けている。厳密同一点での比較ではない。
      (\code{compare\_experiment\_*\_vs\_*P.py} は同一格子の CSV を
      前提に重ね描きする。)
\item \textbf{\code{--alpha-gamma-map} の既定対応}
      ($\alpha=0.1\!\to\!\gamma=1$, $\alpha=1\!\to\!100$, $\alpha=10\!\to\!1000$)
      は $\gamma$ を $\alpha$ とともに大きくする設定であり,
      実験 2-P / 3-P の \code{\_gmap} 系 CSV・図はこの設定で生成されている。
      一律 $\gamma=5.0$ の結果 (\code{\_g5.0}) とは別物。
\item \textbf{ドキュメントの軽微な不整合}:
      \code{experiment\_2P\_delayoff.py} の docstring は
      「$\beta$ 対数走査 40 点」と書いているが, 実装の既定は 15 点
      (\code{BETA\_N\_POINTS = 15}) であり README も 15 点。
\end{enumerate}

% =====================================================================
\section*{参考文献 (実装が参照しているもの)}
% =====================================================================
\begin{itemize}
\item O'Cinneide, C. A. (1993). Entrywise perturbation theory and error analysis
      for Markov chains.
\item O'Cinneide, C. A. (1996). Relative-error bounds for the LU decomposition
      via the GTH algorithm.
\item Fischer, W., Meier-Hellstern, K. (1993). The Markov-modulated Poisson
      process (MMPP) cookbook.
\item Le-Anh, T., Phung-Duc, T. (2025). (コスト関数および実験パラメータ規模の出典)
\end{itemize}

\end{document}
